\documentclass[10pt]{article}
\usepackage{graphicx}
\usepackage{amsmath}
\usepackage{amssymb}
\usepackage{amsfonts}
\usepackage{textcomp}
\usepackage{amsthm}
\usepackage{mathtools}

\theoremstyle{plain}
\newtheorem{theorem}{Theorem}[section]
\newtheorem{lemma}[theorem]{Lemma}
\newtheorem{corollary}[theorem]{Corollary}
\newtheorem{proposition}[theorem]{Proposition}

\theoremstyle{definition}
\newtheorem{definition}[theorem]{Definition}
\newtheorem{example}[theorem]{Example}

\theoremstyle{remark}
\newtheorem{remark}[theorem]{Remark}
\title{%
	\vspace{-1.5em}%
	\large\textbf{Additional Definitions}%
}

\author{A. E. Mahrouss\\amlal@nekernel.org}
\date{\today}

\begin{document}
	
\maketitle
\vspace{-0.5em}
\vspace{0.8em}
\rule{\linewidth}{0.4pt}
\vspace{1em}
	
	\section{Equations \& Definitions}
	
	\begin{definition}
		For:
		\begin{align*}
			\sqrt{2}+2&=\sqrt{2} + x;\\
			2\sqrt{2}+2 &=x;\\
			2(\sqrt{2} + 1) &= x;
		\end{align*}
		Then. We can solve for $l=x=2(\sqrt{2} + 1).$
		\begin{align*}
			l + r &\leq n - 1;\\
			2(\sqrt{2} + 1) + r &\leq n - 1.
		\end{align*}
		Also. There exists for $r \in \mathbb{D}^{+}.$
		\begin{align*}
			2(\sqrt{2} + 1) + \frac{\frac{n - 1}{2} - 4}{2}  &\leq \frac{n - 1}{2} + n - 1;\\
			2(\sqrt{2} + 1) + \frac{\frac{n - 1}{2} - 4}{2} - \frac{n - 1}{2} &\leq n - 1.
		\end{align*}
	\end{definition}
	
	\nocite{*}
	
	\bibliographystyle{acm}
	\bibliography{A_Course_of_Pure_Mathematics, General_Theory_Dirichlet_Series, Riemann}
	
\end{document}